% SEGMENT OLP-0387-S01
% Part: many-valued-logic
% Chapter: syntax-and-semantics
% Section: formulas

\documentclass[../../../include/open-logic-section]{subfiles}

\begin{document}

\olfileid{mvl}{syn}{fml}

\olsection{\usetoken{P}{formula}}

\begin{defn}[வாய்பாடு]
\ollabel{defn:formulas}
கூற்று மொழி~$\Lang L$ இன் \emph{!!{formula}s} அடங்கிய கணம்~$\Frm[L]$
பின்வருமாறு தொகுத்தறிதலாக வரையறுக்கப்படுகிறது:
\begin{enumerate}
\item ஒவ்வொரு !!{propositional variable}~$\Obj p_i$ உம் ஓர் அணு
  !!{formula}.
\item $\Lang L$ இன் ஒவ்வொரு $0$-இட இணைப்பியும் (கூற்று மாறிலியும்)
  ஓர் அணு !!{formula}.
\item $\star$ என்பது~$\Lang L$ இன் ஓர் $n$-இட இணைப்பியாகவும், $!A_1$,
\dots, $!A_n$ ஆகியவை !!{formula}s ஆகவும் இருந்தால்,
$\star(!A_1, \dots, !A_n)$ என்பது !!a{formula}.
\tagitem{limitClause}{வேறு எதுவும் !!a{formula} அன்று.}{}
\end{enumerate}
$\star$ ஒரு $1$-இட இணைப்பி எனில், $\star(!A_1)$ என்பதைப் பெரும்பாலும்
$\star !A_1$ என்றே எழுதுவோம். $\star$ ஒரு $2$-இட இணைப்பி எனில்,
$\star(!A_1,!A_2)$ என்பதைப் பெரும்பாலும் $(!A_1 \star !A_2)$ என்று
எழுதுவோம்.
\end{defn}

வழக்கம்போல், மிக வெளிப்புற அடைப்புக்குறிகளைப் பல நேரங்களில் எழுதாமல்
விடுவோம்.

\begin{ex}
  செந்தர மொழி~$\Lang{L_0}$ இல், $\Obj p_1 \lif (\Obj p_1
  \land \lnot \Obj p_2)$ ஒரு வாய்பாடு. பெருக்கல் தருக்கத்தின் மொழியில்
  அதை $\Obj p_1 \lif (\Obj p_1 \odot
  \lnot \Obj p_2)$ என்று எழுதுவோம். மொழியில் $1$-இட
  $\triangle$ ஐச் சேர்த்தால், $\triangle (\Obj p_1
  \land \Obj p_2) \lif (\triangle \Obj p_1 \land \triangle \Obj p_2)$
  போன்ற வாய்பாடுகளும் கிடைக்கும்.
\end{ex}

\end{document}
